\section{Security Analysis and Residual Risks}
\label{sec:security}
\paragraph{Forged channel identity.}
An injected or self-signed peer must pass configured certificate validation
and prove key possession before the gateway obtains a verified identity.
Hostname validation prevents accepting a different service under the same
root. Tests exercise missing, expired, untrusted, wrong-purpose, and
key-mismatched credentials. This evidence checks implementation paths; it
does not quantify all possible authentication failures.

\paragraph{Replay and token substitution.}
QUIC packet protection and stream recovery handle protected transport bytes;
the policy additionally checks unique grant IDs and live challenges. An
intercepted grant bound to one certificate cannot simply be presented by an
unrelated certified key. The PAP profile also binds exact arguments and checks
current session permission. The standalone grant is intentionally narrower:
it authorizes a fixed read, not a general signed command. Reusing one key
across installations enlarges the compromise and cross-session attack surface.

\paragraph{Speech and model manipulation.}
An adversarial utterance may still influence either model. The protected
gateway can deny a proposed action outside its grant, provided all effects
are mediated there. It does not make transcripts trustworthy, prevent a
model from revealing data through another tool, or stop an authorized but
harmful in-scope action. A business employee who independently acts on speech
is outside the gateway invariant. The evaluation must test these boundaries
explicitly rather than report rejected schema probes as a general prompt
injection defense.

\paragraph{Relays, call correlation, and physical origin.}
Fresh challenges prevent reuse of a stale control exchange but do not prevent
an online adversary from forwarding a live one. A successful exchange
authenticates the remote key across an audio path; it does not authenticate
the route, speaker, physical location, or ownership of all concurrent speech.
Binding to a trusted call identifier or transcript would require additional
design and a careful relay threat model. Such bindings are not implemented.

\paragraph{Denial of service and codec interference.}
CRC, size checks, partial-packet expiry, queue limits, and deadlines bound
several buffering paths. They do not defeat deliberate jamming or guarantee
fair resource allocation before authentication. Traffic length and timing
can expose activity patterns. Our prototype does not hide that modem traffic
is present, provide traffic-analysis resistance, or certify safety for
hearing aids and other acoustic equipment.

\paragraph{Revocation and commit races.}
The business must consult current grant/certificate state, not just accept a
previously completed TLS handshake. The PAP revalidation hook reduces the
window introduced by asynchronous risk assessment. An atomic relationship
between the last authorization check and an external business effect still
requires the business's transaction design. No protocol can turn an
unacknowledged response into evidence that an operation did not happen.
